
\section{Overview}
\label{sec:overview}

\subsection{Basic Concepts in \ours}
\label{sec:concept}

This section introduces the primary concepts in \ours: {\it message}, {\it agent}, {\it service}, and {\it workflow}.
These four concepts are throughout the platform and all multi-agent applications based on it. 

\begin{itemize}
\item \textbf{Message}:
Messages serve as the carriers for information exchange in multi-agent conversations, encapsulating the source and content of the information. In \ours, messages are implemented as Python dictionaries with two mandatory fields (\textit{name} and \textit{content}) and an optional field (\textit{url}). 
The \textit{name} field records the name of the agent that generates the message, and the \textit{content} field contains the text-based information generated by the agent.
The \textit{url} field is designed to hold the Uniform Resource Locator (URL), which typically links to multi-modal data, such as images or videos.
Messages with this field are particularly relevant for interactions with agents that can process and generate multi-modal content. 
Each message is uniquely identified by an auto-generated UUID and timestamp, ensuring traceability. 
\lstlistingname~\ref{lst:msg} shows how the messages can be created, 
serving as atoms in the inter-agent communication of \ours.

\begin{lstlisting}[language=Python, caption={Illustrative examples of message creation in \ours.}, label={lst:msg}, float=h]
from agentscope.message import Msg

msg1 = Msg("Alice", "Hello!")
msg2 = Msg(
    name="Bob",
    content="How do you find this picture I captured yesterday?",
    url="https://xxx.png"
)
\end{lstlisting}

\item \textbf{Agent}: 
Agents are the primary actors within multi-agent applications, acting as conversational participants and executors of tasks. In \ours, agent behaviors are abstracted through two interfaces: the \textit{reply} and \textit{observe} functions. The \textit{reply} function takes a message as input and produces a response, while the \textit{observe} function processes incoming messages without generating a direct reply. 
The interplay between agents and messages, as shown in \lstlistingname~ \ref{lst:agent}, forms the operational basis of \ours and is essential for developers to model complex interactions in multi-agent LLMs.

\begin{lstlisting}[language=Python, caption={Demonstration of message exchange between agents in \ours.}, label={lst:agent}, float=h]
# agent1 and agent2 are two initialized agents, for example
# agent1, agent2 = DialogAgent(...), DialogAgent(...)
msg1 = agent1()
msg2 = agent2(msg1)
\end{lstlisting}

% \item \textbf{Service}: 
% Services in \ours refer to the functional APIs that enable agents to perform specific tasks. These services are categorized into model API services, which are channels to use the LLMs, and general API services, which provide a variety of \emph{tool} functions. The integration of these services into agents is key for executing a wide range of tasks, especially when interfacing with LLMs that may require external data or computation services.

\item \textbf{Workflow}:
Workflows represent ordered sequences of agent executions and message exchanges between agents, analogous to computational graphs in TensorFlow, but with the flexibility to accommodate non-DAG structures. 
Workflows define the flow of information and task processing among agents, facilitating parallel execution and efficiency improvements. 
This concept is essential for designing multi-agent systems that interact with LLMs, as it allows for the coordination of complex, interdependent tasks.


\item \textbf{Service Functions and Tools}: 
Note that service functions are closely related to but \emph{different} from the concept, tools, in the context of agent design in \ours.
Service functions refer to the functional APIs that return a formatted output \texttt{ServiceResponse}, while tools refer to processed services functions with functionality descriptions and necessary input parameters prepared.
We introduce these two concepts in \ours because LLMs require help to invoke service functions as tools.
One observation is that LLMs may need help understanding the functionalities of the service functions precisely and demand more descriptive information to make accurate decisions.
Meanwhile, LLMs can not (reliably) fill in some input parameters of the APIs, such as the API keys of Bing and Google Search.
As a result, \ours defines tools as processed service functions.
% These services are categorized into model API services, which are channels to use the LLMs, and general API services, which provide a variety of \emph{tool} functions. The integration of these services into agents is key for executing a wide range of tasks, especially when interfacing with LLMs that may require external data or computation services.
\end{itemize}

\begin{figure}
    \centering
    \includegraphics[width=0.6\linewidth]{figures/arch4.pdf}
    \caption{Architecture of \ours.}
    \label{fig:arch}
\end{figure}

\subsection{Architecture of \ours}
We present \ours as an infrastructural platform to facilitate the creation, management, and deployment of multi-agent applications integrated with LLMs. 
The architecture of \ours comprises three hierarchical layers and a set of user interaction interfaces, as shown in \figref{fig:arch}.
These layers provide support for multi-agent applications from different levels, including elementary and advanced functionalities of a single agent (utility layer), resources and runtime management (manager and wrapper layer), and agent-level to workflow-level programming interfaces (agent layer).
\ours introduces intuitive abstractions designed to fulfill the diverse functionalities inherent to each layer and simplify the complicated inter-layer dependencies when building multi-agent systems.
Furthermore, we offer programming interfaces and default mechanisms to strengthen the resilience of multi-agent systems against faults within different layers.

\begin{itemize}
\item \textbf{Utility Layer}: As the platform's foundation, the utility layer in \ours provides essential services to support the core functionalities of agents. This layer abstracts the complexity of underlying operations, such as model API invocation and service functions including code execution and database operations, allowing agents to focus on their primary tasks. 
\ours's utility layer is designed with ease of use and robustness as its utmost priority, supporting versatile operations in multi-agent systems and providing \emph{built-in} autonomous retry mechanisms for exception and error handling against unexpected interruptions.

\item \textbf{Manager and Wrapper Layer}: As an intermediary, the manager and wrapper abstraction layer manages the resources and API services, ensuring high availability of resources and providing resistance to undesired responses from LLMs. 
Unlike the utility layer, which provides default handlers, the manager and wrapper layer also offers customizable interfaces for fault tolerance controls depending on developers' needs and the specific requirements of the application. 
This layer is responsible for maintaining the operational integrity of the agents, a crucial aspect for LLMs to perform consistently under diverse conditions. Detailed elaboration on the fault tolerance mechanisms is provided in Section \ref{sec:tolerance}.

\item \textbf{Agent Layer}: At the core of \ours lies the agent abstraction, which forms the backbone of the multi-agent workflow and is the primary entity responsible for interaction and communication. This layer is designed to facilitate the construction of intricate workflows and enhance usability, reducing the programming burden on developers. 
By integrating streamlined syntax and tools, \ours empowers developers to concentrate on the implementation and optimization of agent-based applications that leverage the capabilities of LLMs. 
The programming features and syntactic sugars are introduced in Section \ref{sec:usability} with more details.

\item \textbf{User interaction}: In addition to the layered architecture, \ours provides multi-agent oriented interfaces such as an annotated terminal presenting basic information, Web UI monitoring the system, a Gradio-base~\citep{abid2019gradio} interface that can change a command line application to a graphical one with only one step and a drag-and-drop zero-code programming workstation (Figure~\ref{fig:workstation}). These interfaces allow developers to effortlessly monitor the status and metrics of the application, including agent communication, execution timing, and financial costs.

\end{itemize}

Collectively, the layered constructs of \ours provide the essential building blocks for developers to craft bespoke multi-agent applications that leverage the advanced capabilities of large language models. The subsequent section will delve into the features of \ours that enhance the programming experience for multi-agent application development.
